\documentclass[11pt]{article}
\usepackage[margin=1in]{geometry}
\usepackage{amsmath,amssymb,amsfonts}
\usepackage{booktabs}
\usepackage{array}
\usepackage{hyperref}
\usepackage{xcolor}
\usepackage{enumitem}
\usepackage{longtable}
\usepackage{microtype}
\hypersetup{colorlinks=true,linkcolor=blue,citecolor=blue,urlcolor=blue}

\title{The W33 HoloNet Q4 Diamond:\\A Machine-Audited Master Synthesis with Rolling-Epoch Repair}
\author{Wil Dahn\\W33 Theory Project}
\date{June 19, 2026}

\newcommand{\Qfour}{Q_4}
\newcommand{\Qthree}{Q_3}
\newcommand{\Ftwo}{\mathbb{F}_2}
\newcommand{\Cech}{\check{C}}
\newcommand{\Hcech}{\check{H}}

\begin{document}
\maketitle

\begin{abstract}
This note gives a professional, machine-audited synthesis of the W33 HoloNet Q4 diamond.  The architectural chain is
\[
  \text{topology} \longrightarrow \text{code} \longrightarrow \text{algebra} \longrightarrow \text{physics},
\]
with the 4-cube \(\Qfour\) driving the claimed parameters of a photonic stabilizer-code architecture.  The note consolidates BT1326, audits its number table with BT1327, repairs the synchronization-epoch derivation with BT1328, and separates certificate-backed arithmetic from structural, simulation, and engineering claims using BT1331.  The main correction is that the master epoch
\[
  10980
\]
is not the literal least common multiple \(\mathrm{lcm}(3660,1620)\), which equals \(98820\).  Instead, it is the three-frame rolling chart-phase closure
\[
  3660=6\cdot 540+180,\qquad 3\cdot180=540,\qquad 3\cdot3660=10980.
\]
The corrected synthesis is stronger because every exact numerical claim is either verified by a machine audit or isolated as a separate simulation or engineering gate.
\end{abstract}

\tableofcontents

\section{Executive Summary}

The W33 HoloNet diamond is the claim that a single combinatorial input, the 4-dimensional hypercube \(\Qfour\) with icosahedral chart symmetry, determines an integrated quantum-error-correcting and photonic architecture.  The slogan is
\[
\boxed{\text{TOPOLOGY}\to\text{CODE}\to\text{ALGEBRA}\to\text{PHYSICS}.}
\]
In BT1326 this was stated as the first master synthesis of the W33 HoloNet.  In this paper we convert that statement into a certificate-oriented object.

The audited diamond has four pillars:
\begin{enumerate}[label=(\roman*)]
  \item \textbf{Topology.}  The base object is \(\Qfour\), with \(|V(\Qfour)|=16\) and \(|E(\Qfour)|=32\).  A 540-chart \(\Qthree\) atlas supplies the global chart cover.
  \item \textbf{Code.}  The code layer is the asserted \([[32,4,4]]\) CSS/subsystem code: 32 physical degrees of freedom, 4 logical qubits, and distance 4.
  \item \textbf{Algebra.}  The Clifford restriction \(\mathrm{Cl}(\Qfour)\to\mathrm{Cl}(\Qthree)\) produces an 8-dimensional local spinor object and a global-section logical qubit represented cohomologically.
  \item \textbf{Physics.}  The proposed implementation uses 8 graded waveguide modes per chart, linear optical operations, and a synchronized photonic clock.
\end{enumerate}

The central technical improvement relative to BT1326 is the introduction of a certificate boundary.  Some claims are exact arithmetic.  Others are structural theorem claims already represented in proof notes.  Still others, especially the claimed \(14.4\%\) loss threshold and the \(5\text{mm}\times 5\text{mm}\) chip footprint, are not exact arithmetic and should be treated as simulation and engineering gates.

\section{The Q4 Diamond}

Let \(\Qfour\) be the 4-dimensional hypercube.  Its elementary counts are
\[
|V(\Qfour)|=2^4=16,
\qquad
|E(\Qfour)|=4\cdot 2^3=32.
\]
The synthesis interprets the 32 edges as the physical-qubit budget of the base code block.  The four coordinate directions provide the logical rank used in the W33 code interpretation, while distance 4 is identified with the four-direction hypercube connectivity scale.

The chart layer uses a 540-chart \(\Qthree\) atlas.  The audited arithmetic treats the 540 count as
\[
540 = 120\cdot \frac{9}{2},
\]
where 120 is the order of the icosahedral reflection group \(W(H_3)\).  In the repository's theorem chain this count is also related to the 2160-slot mirror bus by
\[
540 = \frac{2160}{4}.
\]

\section{Machine-Audited Number Table}

BT1327 audits the exact arithmetic portion of the BT1326 number table.  Table~\ref{tab:number-table} records the values that are directly arithmetic-checkable.

\begin{longtable}{>{\raggedright\arraybackslash}p{0.22\linewidth} >{\raggedright\arraybackslash}p{0.25\linewidth} >{\raggedright\arraybackslash}p{0.38\linewidth}}
\caption{Machine-audited Q4 diamond numerical identities.}\label{tab:number-table}\\
\toprule
Number & Role & Audited derivation\\
\midrule
\endfirsthead
\toprule
Number & Role & Audited derivation\\
\midrule
\endhead
16 & Router states & \(|V(\Qfour)|=2^4\)\\
32 & Physical qubits & \(|E(\Qfour)|=4\cdot 2^3\)\\
4 & Logical qubits & Q4 coordinate/logical rank in the W33 interpretation\\
4 & Distance & Four-direction hypercube connectivity scale\\
540 & Atlas charts & \(120\cdot 9/2\)\\
2160 & Mirror slots & \(540\cdot 4\) consistency check\\
14641 & Ihara marker & \(11^4\)\\
3240 & Pair intersections & \(540\cdot12/2\)\\
6480 & Syndrome bits & \(3240\cdot2\)\\
1620 & Independent syndromes & \(6480/4\)\\
268 & Logical qubits/revolution & \(67\cdot4\)\\
8 & Spinor modes/chart & \(2^3\)\\
4320 & Waveguide channels & \(540\cdot8\)\\
70,778,880 & Three-level modes & \(540\cdot4\cdot32^3\)\\
\bottomrule
\end{longtable}

The audit result is intentionally strict.  It does not silently pass an incorrect derivation.  The only failing literal check in BT1327 was the statement
\[
10980=\mathrm{lcm}(3660,1620).
\]
A direct calculation gives
\[
\mathrm{lcm}(3660,1620)=98820.
\]
Thus the epoch value could not be justified by the LCM line as written.

\section{Rolling-Epoch Repair}

BT1328 restores the value 10980 with the correct mechanism.  The relevant observation, already latent in BT1321, is that an Ihara frame of 3660 logical-qubit frames does not align integrally with the 540-chart atlas.  Instead,
\[
3660=6\cdot540+180,
\]
so each Ihara frame advances the chart atlas by an offset of 180 chart slots.  Since
\[
180=\frac{540}{3},
\]
three such offsets close one full chart cycle:
\[
3\cdot180=540.
\]
Therefore the synchronization epoch is
\[
3\cdot3660=10980.
\]
This is a rolling phase closure, not an LCM closure.  The corrected theorem statement is:
\begin{quote}
The master synchronization epoch is 10,980 Ihara sub-periods, derived by three-frame rolling chart-phase closure: \(3660=6\cdot540+180\), \(3\cdot180=540\), and \(3\cdot3660=10980\).
\end{quote}

\section{Cohomological Logical Qubit}

The global logical qubit is represented in the proof chain as a non-localizable section of the chart atlas.  Let \(S\) denote the spinor bundle obtained from the Clifford restriction
\[
\mathrm{Cl}(\Qfour)\to\mathrm{Cl}(\Qthree).
\]
BT1323 identifies the global logical operator \(\gamma_4\) with a nonzero cohomology class
\[
[\gamma_4]\in \Hcech^1(\{U_\alpha\},S).
\]
The 540-chart atlas has 12 chart-neighbors per chart, hence
\[
\#\{\text{pair intersections}\}=\frac{540\cdot12}{2}=3240.
\]
With a 2-dimensional local spinor contribution on each pair intersection, the Cech 1-cochain dimension is
\[
\dim \Cech^1=3240\cdot2=6480.
\]
This exactly matches the 6480 syndrome-bit count.  Dividing by the distance-4 redundancy scale gives
\[
6480/4=1620
\]
independent syndrome bits.

\section{Photonic Mode Interpretation}

BT1324 maps the local spinor object to 8 waveguide modes per chart.  The basic arithmetic is
\[
\dim \mathrm{Cl}(\Qthree)=2^3=8,
\]
so a 540-chart implementation requires
\[
540\cdot8=4320
\]
base waveguide channels before concatenation overhead.  At three concatenation levels, the stated mode budget is
\[
540\cdot4\cdot32^3=70,778,880,
\]
which is approximately 70.8 million modes.

This paper does not treat the 5mm by 5mm chip claim as exact arithmetic.  That claim requires a physical layout model: waveguide pitch, bend radius, loss, crossing penalty, crosstalk, thermal phase-shifter footprint, routing overhead, and packaging constraints.  It is therefore an engineering gate, not a theorem of \(\Qfour\) alone.

\section{Certificate Boundary}

BT1331 classifies the master synthesis claims into four categories:
\begin{enumerate}[label=(\arabic*)]
  \item \textbf{Exact arithmetic.}  These claims are machine-audited by BT1327 and BT1328.
  \item \textbf{Structural theorem.}  These claims rely on proof notes such as the spinor cohomology and Clifford restriction documents.
  \item \textbf{Simulation gate.}  The \(14.4\%\) threshold is a decoder/noise-model claim and requires a decoder certificate.
  \item \textbf{Engineering gate.}  The chip footprint claim requires layout and foundry feasibility analysis.
\end{enumerate}

This classification makes the master theorem harder to misuse.  Exact counts remain exact.  Structural identifications remain proof-theoretic.  Threshold and fabrication claims become explicit gates.

\section{Implications for the Next Breakthroughs}

The corrected synthesis points to three immediate next steps:
\begin{itemize}
  \item Build a decoder threshold certificate.  A Gottesman-Knill decoder can efficiently simulate stabilizer dynamics, but it cannot violate information-theoretic capacity bounds for photon loss treated as erasure.
  \item Build a foundry feasibility gate.  The 4320 base channels and 70.8M concatenated modes must be compared to pitch, bend radius, loss, crosstalk, and thermal footprint.
  \item Promote the audited paper into the repository build system so the technical note is regenerated by CI.
\end{itemize}

\section{Conclusion}

The W33 HoloNet Q4 diamond is now stronger than the initial master synthesis.  BT1327 identified the one false literal arithmetic derivation.  BT1328 repaired the epoch while preserving the value 10980.  BT1330 patched the source theorem statements.  BT1331 classified the master theorem into exact, structural, simulation, and engineering layers.

The resulting message is precise:
\[
\boxed{10980=3\cdot3660\text{ by rolling chart-phase closure, not }\mathrm{lcm}(3660,1620).}
\]
The exact arithmetic core of the Q4 diamond is machine-audited, and the remaining physical-performance claims are cleanly exposed as certifiable next gates.

\end{document}
